\subsection*{A local context gate supplies useful spatial credit in conductance trees}

The preceding comparisons establish the value of spatially selective credit when its coefficients are supplied. We next asked whether a simple local cue could provide the required distinction in a conductance tree. Four terminal compartments fed two proximal compartments and a soma. Both subtrees received the same positive feature pairs, $\exp(\pm z_j)$, for two independent latent inputs. Binary context increased inhibition in one proximal compartment. Teacher E/I tuning produced increasing responses in both subtrees for an aligned task, or increasing responses in one and decreasing responses in the other for an opposed task (Fig.~\ref{fig:conductancecredit}A,B). Teacher and student shared the circuit, ensuring representability.

The local rule combined a somatic error $\delta$ with the inhibitory activity $a_{p(j)}$ at the parent of terminal compartment $j$:
\begin{equation}
\widehat\delta_j^V=\delta\,\mathbf1[a_{p(j)}=0].
\label{eq:localconductancegate}
\end{equation}
Both proximal compartments and the soma retained the common error. Thus only distal credit was selected by context; all 24 conductance parameters remained plastic. Local eligibilities were identical across rules. This update required the imposed inhibitory cue and somatic error, but neither exact path derivatives nor projection of an exact field during learning.

After fixing the rule and controls, we tested twenty new paired seed blocks at a common Adam rate. Within 4,096 updates, mean opposed-task NMSE was $2.33\times10^{-5}$ with the local gate, $2.18\times10^{-5}$ with exact credit, and 0.487 with uniform broadcast (Fig.~\ref{fig:conductancecredit}C,D). The gate-minus-exact difference was $1.49\times10^{-6}$ (95\% paired interval, $-1.60\times10^{-6}$ to $4.86\times10^{-6}$). On aligned targets, their NMSEs were $1.06\times10^{-5}$, $7.53\times10^{-6}$ and $5.53\times10^{-4}$, respectively. The opposed-minus-aligned difference in the broadcast-minus-gate gap was 0.487 (0.360--0.615), positive in every seed (Fig.~\ref{fig:conductancecredit}E). Continuing all conditions to 16,384 updates retained an interaction of 0.468 (0.332--0.604). Two additional frozen rates gave the same qualitative result (Supplementary Fig.~\ref{fig:supp_local_gate_rates}).

The controls identify which distinction mattered. Reversing the distal selector raised opposed-task NMSE to 0.949. Applying the selector to proximal credit as well gave 0.123 and froze both inhibitory gains: their eligibility requires inhibitory activity, precisely when that gate is closed (Fig.~\ref{fig:conductancecredit}F). Two oracle-controlled leaf profiles with common proximal credit achieved $2.28\times10^{-5}$, showing that a separate projected proximal pattern was unnecessary. A continuous gate based on local inhibitory conductance also learned accurately. These comparisons support a supplied branch selector in this circuit, not autonomous generation of a biological teaching signal or a general minimum feedback rank. Exact and fixed-profile historical controls, including the first monotonic construction, remain in the supplement.
